% Acronyms
\newacronym{api}{API}{Application Program Interface}
\newacronym{uml}{UML}{Unified Modeling Language}
\newacronym{tsa}{TSA}{Termine solo acronimo}
\newacronym{wcag}{WCAG}{Web Content Accessibility Guidelines}
\newacronym{agid}{AGID}{Agenzia per l'Italia Digitale}
\newacronym{llm}{LLM}{Large Language Model}
\newacronym{eaa}{EAA}{European Accessibility Act}
\newacronym{ict}{ICT}{Information and Communication Technology}
\newacronym{w3c}{W3C}{World Wide Web Consortium}
\newacronym{pour}{POUR}{Percepibile, Operabile, Comprensibile, Robusto}
\newacronym{iso}{ISO}{International Organization for Standardization}
\newacronym{ue}{UE}{Unione Europea}
\newacronym{uni}{UNI}{Ente Italiano di Normazione}
\newacronym{cei}{CEI}{Comitato Elettrotecnico Italiano}
\newacronym{dlgs}{D.Lgs.}{Decreto Legislativo}
\newacronym{wave}{WAVE}{Web Accessibility Evaluation Tool}
\newacronym{mauve}{MAUVE}{Multiguideline Accessibility and Usability Validation Environment}

% Glossary
\newglossaryentry{tranco}{
    name={Ranking Tranco},
    text={ranking Tranco},
    description={Classifica dei domini Internet più popolari, calcolata combinando i dati di più fonti (anziché di una sola, come nelle classifiche tradizionali del tipo Alexa) allo scopo di ridurre la volatilità giornaliera e la possibilità di manipolazione artificiale della graduatoria; utilizzata, ad esempio, dal WebAIM Million come base per selezionare il campione di siti da analizzare}
}
\newglossaryentry{agid-descrizione}{
    name={AGID},
    text={AGID},
    description={Agenzia per l'Italia Digitale, ente pubblico italiano responsabile, fra le altre attività, del monitoraggio periodico sull'accessibilità dei siti web e delle applicazioni mobili delle pubbliche amministrazioni, ai sensi della legge Stanca e del D.Lgs.~82/2022, e della pubblicazione delle relative Linee guida}
}
\newglossaryentry{wave-descrizione}{
    name={WAVE},
    text={WAVE},
    description={Strumento di valutazione automatica dell'accessibilità sviluppato da WebAIM, utilizzato per analizzare le home page nell'ambito del report WebAIM Million citato nel presente lavoro}
}
\newglossaryentry{mauve-descrizione}{
    name={MAUVE},
    text={MAUVE},
    description={Strumento di validazione automatica dell'accessibilità di pagine web e documenti digitali (in particolare PDF) rispetto alle WCAG, sviluppato dal laboratorio HIIS del CNR-ISTI e utilizzato da AGID nell'ambito del monitoraggio nazionale sull'accessibilità dei siti delle pubbliche amministrazioni}
}
\newglossaryentry{disciplinaorganica}{
    name={disciplina organica},
    text={disciplina organica},
    description={Corpo normativo che regola una materia in modo sistematico e complessivo, individuando principi generali e relative modalità di attuazione, anziché intervenire attraverso disposizioni frammentarie o settoriali; espressione riferita, nel presente lavoro, alla legge Stanca in quanto prima normativa italiana a disciplinare in modo organico l'accessibilità degli strumenti informatici}
}